% explainer-source-hash: e86e98979cf8c7389d590088e201dc0e7817163aa2a1f304cecfbf7c32a770ce
% explainer-source-commit: 3f1d18ace6744e68643813b24e0dc406bd4e6c7c
% explainer-generated: 2026-07-16
% Regenerate with /explain-all (see WIRING.md). Do not edit this stamp by hand:
% it is how the toolchain knows whether this document still matches the code.
\documentclass[11pt,a4paper]{article}
\input{preamble}

\title{\textbf{Voice Appointment Agent}\\[0.3em]
\large A Deep Dive: how a phone call becomes a booked appointment}
\author{Portfolio explainer, built from the repository source}
\date{}

\begin{document}
\maketitle
\thispagestyle{empty}

\begin{keyidea}{The one-sentence version}
Somebody phones a small business. Twilio answers the call and turns the caller's speech
into text. That text is handed to Google's Gemini 2.5 Flash language model, which works out
what the caller actually wants (book, reschedule, or cancel) and pulls out the details (name,
date, time). A small database then really does book the appointment, refusing to
double-book. The whole thing is designed so that the expensive parts (the phone line and the
language model) can be unplugged and replaced with free stand-ins, which is why it can be
tested 54 times over without spending money.
\end{keyidea}

\tableofcontents
\newpage

\section{What this project is, and what it is not}

\subsection{The problem it addresses}

A clinic, a salon, or a repair shop loses bookings when nobody picks up the phone. The
caller does not leave a voicemail; the caller phones the next business on the list. The
traditional fix is a ``phone tree'' (``press 1 for bookings''), which callers dislike and
which cannot handle a sentence like ``can you move my Tuesday thing to Thursday morning?''.

This project is a working demonstration of the modern alternative: a real telephone number
that a real person can speak to in ordinary sentences, backed by a language model that
understands the request, and by a calendar that genuinely reserves the slot.

\subsection{Scope, stated honestly and up front}

\begin{honest}{This is a demonstration of a pattern, not a deployed product}
Read the repository carefully and the boundary is clear, and the project's own README is
unusually candid about it. What is genuinely proven: the conversation logic, the calendar,
the conflict detection, the HTTP endpoints, and that the real Gemini model classifies clean
typed English correctly. What is \emph{not} proven: that any of this survives contact with a
real human voice on a real phone line. Twilio's speech-to-text output is noisy (mumbling,
background noise, half-sentences), and the language model in this repository has only ever
been shown clean, hand-typed sentences. That gap, between ``the API call works'' and ``this
works on a phone call'', is real and the repository says so itself.
\end{honest}

\begin{jargon}{Telephony}
The technology of telephone calls. Here it means the machinery that makes a physical phone
ring and carries the audio, which this project rents from a company called Twilio rather than
building itself.
\end{jargon}

\subsection{Size and shape of the codebase}

Roughly 1,900 lines of Python across nine source modules, five test files, and a
25-scenario evaluation harness. It is a compact, focused project. Every tracked file was
read in full to build this document.

\begin{figure}[H]
\centering
\begin{forest} dirtree
[voice-agent
  [voice\_agent/ \quad {\rmfamily\itshape the application itself}
    [app.py \quad {\rmfamily\itshape FastAPI web service Twilio talks to}]
    [conversation.py \quad {\rmfamily\itshape the state machine: all the business logic}]
    [intent\_engine.py \quad {\rmfamily\itshape the Gemini interface + implementation}]
    [calendar\_store.py \quad {\rmfamily\itshape SQLite + Google Calendar backends}]
    [session\_store.py \quad {\rmfamily\itshape per-call memory and transcript}]
    [twiml.py \quad {\rmfamily\itshape builds Twilio's XML instructions}]
    [config.py \quad {\rmfamily\itshape environment variables, business rules}]
    [cli.py \quad {\rmfamily\itshape places a real outbound test call}]
  ]
  [eval/ \quad {\rmfamily\itshape the offline evaluation harness}
    [scenarios.py \quad {\rmfamily\itshape the 25 scripted conversations}]
    [scripted\_engine.py \quad {\rmfamily\itshape free stand-in for Gemini}]
    [run\_eval.py \quad {\rmfamily\itshape prints the pass/fail report}]
  ]
  [tests/ \quad {\rmfamily\itshape 54 pytest tests}
    [test\_calendar\_store.py \quad {\rmfamily\itshape 13 tests, no network}]
    [test\_twiml\_endpoints.py \quad {\rmfamily\itshape 7 tests via a fake HTTP client}]
    [test\_eval\_scenarios.py \quad {\rmfamily\itshape runs all 25 scenarios}]
    [test\_llm\_intent.py \quad {\rmfamily\itshape 6 REAL, billed Gemini calls}]
    [test\_twilio\_integration.py \quad {\rmfamily\itshape signature check + 1 free lookup}]
  ]
  [.env.example \quad {\rmfamily\itshape placeholders only; real secrets live outside}]
  [requirements.txt]
  [README.md]
]
\end{forest}
\caption{The repository. Note the deliberate separation: \file{voice\_agent/} is production
code, \file{eval/} contains a test double that must never run in production.}
\end{figure}

\section{The concepts you need first}

This section defines everything the rest of the document assumes. If a term is already
familiar, skip it.

\begin{jargon}{API (Application Programming Interface)}
An agreed way for one program to ask another program to do something. When this project
``calls the Gemini API'', it sends a message over the internet to Google's servers and gets
a reply back. Google charges a small amount per call.
\end{jargon}

\begin{jargon}{Webhook}
The reverse of the usual arrangement. Normally your program calls someone else's server.
With a webhook, you give \emph{them} a web address, and \emph{they} call \emph{you} when
something happens. Twilio does exactly this: when a call comes in, Twilio makes an HTTP
request to this project's web address and asks ``what should I say?''. This is why the
project is a web server at all.
\end{jargon}

\begin{jargon}{HTTP request / POST}
The message format the web runs on. A ``POST'' request is one that carries a body of data,
as opposed to just asking for a page. Twilio POSTs the details of the call (who is
calling, what they said) to this service.
\end{jargon}

\begin{jargon}{TwiML (Twilio Markup Language)}
Twilio's instruction format: a small XML document telling Twilio what to do next on the
call. \code{<Say>} means speak this text aloud. \code{<Gather>} means listen to the caller
and send me what you heard. \code{<Hangup>} means end the call. This project's entire job,
from Twilio's point of view, is to answer with TwiML.
\end{jargon}

\begin{jargon}{XML}
A text format that stores structured data using angle-bracket tags, like
\code{<Say>Hello</Say>}. TwiML is one specific flavour of it.
\end{jargon}

\begin{jargon}{LLM (Large Language Model)}
A program trained on enormous quantities of text that can read a sentence and produce a
sensible reply. Gemini 2.5 Flash is Google's fast, cheap model. ``Flash'' signals that it is
tuned for speed over maximum capability, which matters here because a caller on a phone will
not tolerate long silences.
\end{jargon}

\begin{jargon}{Slot filling}
The task of collecting a fixed set of facts through conversation. The ``slots'' here are
name, phone, date, time, and duration. The agent's job each turn is to fill whichever slots
are still empty, one question at a time, until it has enough to book.
\end{jargon}

\begin{jargon}{State machine}
A program that is always in exactly one of a limited number of situations (``states''), and
moves between them according to fixed rules when something happens. A conversation is a
natural fit: ``no idea what they want yet'' becomes ``booking, still need a name'' becomes
``booking, waiting for them to confirm''.
\end{jargon}

\begin{jargon}{SQLite}
A complete database that lives in a single ordinary file on disk, with no separate server
process to install or run. Ideal for a small project. Its main limitation, which matters
later, is that only one thing may write to it at a time.
\end{jargon}

\begin{jargon}{Interface (in programming)}
A published list of operations that something promises to provide, without saying how. If
you write your program against the interface ``a calendar can book, cancel and reschedule'',
you can later swap the real calendar for a fake one and your program cannot tell. This single
idea is the backbone of this project's architecture, and Section~\ref{sec:seams} is devoted
to it.
\end{jargon}

\begin{jargon}{Test double}
A stand-in used during testing in place of the real, slow, or expensive thing, the way a
stunt double stands in for an actor. This project has one: a fake ``language model'' that is
really just pattern-matching rules.
\end{jargon}

\begin{jargon}{E.164}
The international standard for writing phone numbers unambiguously: a plus sign, a country
code, then the number, with no spaces or dashes. \code{+15551234567}. Remember this one;
an inconsistency around it is a real bug found in Section~\ref{sec:findings}.
\end{jargon}

\section{Architecture: the shape of the whole system}

\begin{figure}[H]
\centering
\resizebox{\textwidth}{!}{
\begin{tikzpicture}[node distance=9mm]
  \node[extbox] (caller) {Human caller\\{\scriptsize on a real phone}};
  \node[extbox, right=26mm of caller] (twilio) {\textbf{Twilio}\\{\scriptsize answers the call,}\\{\scriptsize speech $\rightarrow$ text,}\\{\scriptsize text $\rightarrow$ speech}};
  \node[box, right=26mm of twilio] (app) {\textbf{app.py}\\{\scriptsize FastAPI web service}\\{\scriptsize /voice/incoming}\\{\scriptsize /voice/gather}\\{\scriptsize /voice/status}};
  \node[box, right=20mm of app] (twiml) {\textbf{twiml.py}\\{\scriptsize builds the}\\{\scriptsize XML reply}};
  \node[box, below=16mm of app] (conv) {\textbf{conversation.py}\\{\scriptsize ConversationManager}\\{\scriptsize ALL business logic}};
  \node[store, right=20mm of conv] (sess) {\textbf{session\_store.py}\\{\scriptsize per-call state}\\{\scriptsize + transcript}};
  \node[box, below left=15mm and 4mm of conv] (intent) {\textbf{intent\_engine.py}\\{\scriptsize IntentEngine}\\{\scriptsize \itshape interface}};
  \node[box, below right=15mm and 4mm of conv] (cal) {\textbf{calendar\_store.py}\\{\scriptsize CalendarStore}\\{\scriptsize \itshape interface}};
  \node[extbox, below=12mm of intent] (gemini) {\textbf{Gemini 2.5 Flash}\\{\scriptsize via Vertex AI}\\{\scriptsize (billed per call)}};
  \node[store, below=12mm of cal] (sqlite) {\textbf{SQLite}\\{\scriptsize appointments}};

  % forward path, labels below the line
  \draw[flow] (caller) -- node[lbl,below]{voice} (twilio);
  \draw[flow] (twilio) -- node[lbl,below]{HTTP POST\\{\scriptsize (a webhook)}} (app);
  \draw[flow] (app) -- node[lbl,right]{one caller\\utterance} (conv);
  \draw[flow] (conv) -- (intent);
  \draw[flow] (conv) -- (cal);
  \draw[flow] (intent) -- node[lbl,right]{JSON} (gemini);
  \draw[flow] (cal) -- (sqlite);
  \draw[flow] (conv) -- node[lbl,above]{save} (sess);
  % return path, arcs above the row so nothing overlaps
  \draw[dflow] (app) to[bend right=24] node[lbl,above]{TwiML response} (twilio);
  \draw[dflow] (twilio) to[bend right=24] node[lbl,above]{spoken reply} (caller);
  % app <-> twiml, a single two-way exchange
  \draw[{Stealth[length=2.4mm]}-{Stealth[length=2.4mm]}, draw=accent, thick, dashed]
    (app) -- node[lbl,above]{reply text\\{\scriptsize $\rightarrow$ TwiML XML}} (twiml);
\end{tikzpicture}}
\caption{The whole system. Solid arrows are requests going in; dashed arrows are responses
coming back. The two boxes marked \emph{interface} are the load-bearing design decision:
everything below them can be swapped out.}
\label{fig:arch}
\end{figure}

\subsection{The single most important structural fact}

\code{conversation.py} contains all the business logic, and it does not know that Twilio or
Gemini exist. It talks only to the two interfaces. Read its import list and this is provable:
it imports \code{IntentEngine} and \code{CalendarStore}, never \code{GeminiIntentEngine} or
any Twilio module.

That is what makes the rest of the project possible, and it is the thing to be able to explain
under questioning. It is unpacked properly in Section~\ref{sec:seams}.

\section{The call flow: what actually happens, in order}

The counter-intuitive part of telephony webhooks is that a single phone call is not one
continuous conversation from the server's point of view. It is a series of \emph{separate,
unrelated HTTP requests}. The server has no memory between them. This drives the whole design
of \code{session\_store.py}.

\begin{keyidea}{Why the server needs a memory at all}
Each time the caller speaks, Twilio makes a brand new HTTP request. As far as the web service
is concerned, that request could be from anyone, about anything. Everything the agent knows
about this call, the name collected two turns ago, the date, whether they confirmed, must
therefore be written to a database and read back on every single turn. The key used to look it
back up is Twilio's \code{CallSid}, a unique identifier Twilio assigns to each call and
includes in every request.
\end{keyidea}

\begin{figure}[H]
\centering
\resizebox{0.97\textwidth}{!}{
\begin{tikzpicture}[node distance=6mm]
  % lifelines
  \node (c) at (0,0) {\textbf{Caller}};
  \node (t) at (4,0) {\textbf{Twilio}};
  \node (a) at (8.6,0) {\textbf{app.py}};
  \node (m) at (12.6,0) {\textbf{Manager}};
  \node (g) at (16.2,0) {\textbf{Gemini}};
  \node (d) at (19.6,0) {\textbf{SQLite}};
  \foreach \n in {c,t,a,m,g,d} { \draw[linegrey, thick] (\n) -- ++(0,-13.4); }

  \draw[flow] (0,-1) -- node[lbl,above]{dials the number} (4,-1);
  \draw[flow] (4,-1.7) -- node[lbl,above]{POST /voice/incoming\\{\scriptsize CallSid, From}} (8.6,-1.7);
  \draw[flow] (8.6,-2.5) -- node[lbl,above]{create session} (19.6,-2.5);
  \draw[dflow] (8.6,-3.2) -- node[lbl,above]{TwiML: \code{<Gather><Say>greeting}} (4,-3.2);
  \draw[dflow] (4,-3.9) -- node[lbl,above]{speaks the greeting, listens} (0,-3.9);

  \node[draw=accent, fill=white, rounded corners=2pt, inner sep=4pt, font=\small\itshape,
        text=inkblue, align=center] at (10,-4.7) {one turn, repeated until the call ends};

  \draw[flow] (0,-5.5) -- node[lbl,above]{``book me Tuesday at 2''} (4,-5.5);
  \draw[flow] (4,-6.2) -- node[lbl,above]{POST /voice/gather\\{\scriptsize SpeechResult=text}} (8.6,-6.2);
  \draw[flow] (8.6,-6.9) -- node[lbl,above]{handle\_turn(text)} (12.6,-6.9);
  \draw[flow] (12.6,-7.6) -- node[lbl,above]{load transcript} (19.6,-7.6);
  \draw[flow] (12.6,-8.3) -- node[lbl,above]{parse\_turn(context)} (16.2,-8.3);
  \draw[dflow] (16.2,-9.0) -- node[lbl,above]{JSON: intent + slots + reply} (12.6,-9.0);
  \draw[flow] (12.6,-9.7) -- node[lbl,above]{book() if all slots filled \& confirmed} (19.6,-9.7);
  \draw[dflow] (12.6,-10.4) -- node[lbl,above]{TurnOutcome(reply, end\_call)} (8.6,-10.4);
  \draw[dflow] (8.6,-11.1) -- node[lbl,above]{TwiML: \code{<Gather>} or \code{<Hangup>}} (4,-11.1);
  \draw[dflow] (4,-11.8) -- node[lbl,above]{speaks the reply} (0,-11.8);

  \draw[flow] (4,-12.9) -- node[lbl,above]{POST /voice/status (call ended)} (8.6,-12.9);
\end{tikzpicture}}
\caption{One call, turn by turn. Every horizontal arrow from Twilio to \file{app.py} is a
\emph{separate HTTP request} with no memory of the last one, which is why the transcript is
re-loaded from SQLite on every turn.}
\label{fig:seq}
\end{figure}

\subsection{The silence problem, and the fix}

This is one of the genuinely non-obvious engineering details in the project, and it is worth
understanding because it is the kind of thing that only surfaces when you read the vendor
documentation properly.

\begin{keyidea}{\code{<Gather>} does not call you back when the caller says nothing}
By default, Twilio's \code{<Gather>} requests its \code{action} URL only when it actually
captures speech. If the caller is silent, Twilio simply gives up and moves on to the
\emph{next} instruction in the TwiML document instead. The consequence is severe: all the
careful ``sorry, could you repeat that?'' retry logic in \code{conversation.py} would be
unreachable for a silent caller, which is precisely the caller who most needs it.

The fix is one attribute, \code{action\_on\_empty\_result=True}, set on every \code{<Gather>}
in \file{voice\_agent/twiml.py}. With it, silence also posts back to \code{/voice/gather},
just with an empty \code{SpeechResult}, so every possible outcome flows through the one state
machine.
\end{keyidea}

\file{twiml.py} is only 40 lines and builds exactly two documents:

\begin{lstlisting}[language=Python,caption={voice\_agent/twiml.py, the essential part}]
def build_gather_twiml(say_text: str, gather_action_url: str) -> str:
    response = VoiceResponse()
    gather = Gather(
        input="speech",
        action=gather_action_url,
        method="POST",
        speech_timeout="auto",
        timeout=GATHER_TIMEOUT_SECONDS,      # 6 seconds
        action_on_empty_result=True,          # <- the fix described above
    )
    gather.say(say_text)
    response.append(gather)
    # Belt and braces: only reached if Twilio somehow does not re-request
    # the action URL (e.g. the call already ended).
    response.say("We didn't receive any input. Goodbye.")
    response.hangup()
    return str(response)
\end{lstlisting}

The trailing \code{<Say>}/\code{<Hangup>} pair is a deliberate safety net. If Twilio ever
does fall through, the call ends politely rather than hanging in silence.

\section{The seams: why two interfaces change everything}
\label{sec:seams}

\subsection{The problem being solved}

The evaluation harness runs 25 scenarios. Several are five turns or more. Run those against
the real Gemini model and each run is over 100 billed API calls, takes real time, and, because
language models are not perfectly deterministic, might give a different answer tomorrow. Doing
that on every code change is unaffordable and unreliable.

\begin{jargon}{Deterministic}
Always producing exactly the same output for the same input. Ordinary code is deterministic;
a language model largely is not, which is why a test that depends on one is fragile.
\end{jargon}

\subsection{The solution}

Define what the conversation manager \emph{needs} rather than what it will \emph{use}:

\begin{lstlisting}[language=Python,caption={The two interfaces, reduced to their essence}]
class IntentEngine(abc.ABC):
    @abc.abstractmethod
    def parse_turn(self, ctx: TurnContext) -> TurnResult: ...

class CalendarStore(abc.ABC):
    @abc.abstractmethod
    def find_conflict(self, start, duration_minutes, exclude_id=None): ...
    @abc.abstractmethod
    def book(self, name, phone, start, duration_minutes) -> Appointment: ...
    @abc.abstractmethod
    def get(self, appointment_id) -> Appointment: ...
    @abc.abstractmethod
    def list_by_phone(self, phone, active_only=True) -> list[Appointment]: ...
    @abc.abstractmethod
    def reschedule(self, appointment_id, new_start, new_duration_minutes=None): ...
    @abc.abstractmethod
    def cancel(self, appointment_id) -> Appointment: ...
\end{lstlisting}

\begin{jargon}{Abstract base class / \code{abc.ABC}}
Python's way of writing an interface. Marking a method \code{@abstractmethod} means ``every
real implementation must provide this''; Python refuses to create an object from a class that
has not filled them all in. It turns a design intention into an error the language enforces.
\end{jargon}

Now the identical \code{ConversationManager} runs in three different worlds:

\begin{figure}[H]
\centering
\resizebox{\textwidth}{!}{
\begin{tikzpicture}[node distance=8mm]
  \node[box, minimum width=42mm] (cm) {\textbf{ConversationManager}\\{\scriptsize the identical code in all three columns}};

  \node[box, below left=14mm and 30mm of cm] (p) {\textbf{1. Production}\\{\scriptsize uvicorn / real calls}};
  \node[box, below=14mm of cm] (t) {\textbf{2. pytest}\\{\scriptsize TestClient, 47 free tests}};
  \node[box, below right=14mm and 30mm of cm] (e) {\textbf{3. eval harness}\\{\scriptsize 25 scenarios}};

  \node[extbox, below=9mm of p, align=center] (p2) {GeminiIntentEngine\\{\scriptsize\itshape real, billed}\\[2pt] SqliteCalendarStore\\{\scriptsize\itshape real file on disk}};
  \node[extbox, below=9mm of t, align=center] (t2) {ScriptedIntentEngine\\{\scriptsize\itshape free regex double}\\[2pt] SqliteCalendarStore\\{\scriptsize\itshape temp file}};
  \node[extbox, below=9mm of e, align=center] (e2) {ScriptedIntentEngine\\{\scriptsize\itshape free regex double}\\[2pt] SqliteCalendarStore\\{\scriptsize\itshape temp file}};

  \draw[flow] (cm) -- (p); \draw[flow] (cm) -- (t); \draw[flow] (cm) -- (e);
  \draw[flow] (p) -- (p2); \draw[flow] (t) -- (t2); \draw[flow] (e) -- (e2);

  \node[lbl, below=5mm of p2, text width=34mm] {costs real money,\\proves the real thing};
  \node[lbl, below=5mm of t2, text width=34mm] {free, instant,\\fully deterministic};
  \node[lbl, below=5mm of e2, text width=34mm] {free, instant,\\fully deterministic};
\end{tikzpicture}}
\caption{One state machine, three wirings. This is dependency injection: the manager is
\emph{handed} its collaborators rather than creating them, so the caller decides whether they
are real or fake.}
\end{figure}

\begin{jargon}{Dependency injection}
A long name for a simple habit. Instead of a class building the things it needs
(\code{self.engine = GeminiIntentEngine()}, which welds it to Gemini forever), it accepts them
as arguments (\code{def \_\_init\_\_(self, intent\_engine, ...)}). Whoever creates the object
chooses what to pass. That choice point is the ``seam'' that testing needs.
\end{jargon}

\begin{honest}{The test double is honest about being one, and that matters}
It would be easy, and dishonest, to present \code{ScriptedIntentEngine} as a ``fallback'' so
the project could claim it works without an API key. The repository refuses to do this. Its
docstring states plainly: \emph{``This is a TEST DOUBLE, not a production fallback.''} It
lives in \file{eval/}, not in \file{voice\_agent/}, so it is physically outside the production
package. \code{app.py} always wires up the real Gemini engine unless a test explicitly passes
something else. The ``LLM-driven'' claim therefore survives scrutiny: the shipped path really
does use the model.
\end{honest}

\section{Module by module}

\subsection{\file{voice\_agent/config.py}}

The smallest module, and quietly disciplined. Every value is read from the environment
\emph{at call time}, never cached at import time, and never logged.

\begin{lstlisting}[language=Python,caption={The required-variable helper and the business rules}]
def _require(name: str) -> str:
    value = os.environ.get(name, "").strip()
    if not value:
        raise ConfigError(f"Required environment variable {name} is not set")
    return value

BUSINESS_HOURS_START = 9   # 9:00 local
BUSINESS_HOURS_END = 17    # 17:00 local
DEFAULT_APPOINTMENT_DURATION_MINUTES = 30
MAX_RETRY_COUNT = 3        # consecutive unclear turns before handoff
MAX_TURN_COUNT = 14        # total turns before forced handoff
\end{lstlisting}

Two small decisions worth noticing. Reading at call time rather than import time is what lets
a test set an environment variable and have it actually take effect. And failing loudly with
\code{ConfigError} when a credential is missing is better than starting up and failing
mysteriously on the first real call.

\begin{jargon}{Environment variable}
A named value supplied to a program by the system that starts it, rather than written in the
code. It is the standard way to supply passwords and API keys, precisely so they never appear
in the source code and never get committed to version control.
\end{jargon}

\subsection{\file{voice\_agent/session\_store.py}}

Two SQLite tables, and the answer to ``how does a stateless web service remember a phone
call?''.

\begin{itemize}
  \item \code{sessions}: one row per call, keyed by \code{call\_sid}. The entire
  \code{SessionState} object is serialised to JSON and stuffed into one \code{state\_json}
  column.
  \item \code{transcript}: one row per utterance, for auditing what was actually said.
\end{itemize}

\begin{jargon}{JSON, and serialisation}
JSON is a plain-text format for structured data. \emph{Serialising} means flattening a
live object into text so it can be stored; \emph{deserialising} is rebuilding the object from
that text. Here, the whole session state becomes one string in one column.
\end{jargon}

\begin{keyidea}{Why store the state as one JSON blob instead of proper columns?}
It is a real trade-off, chosen deliberately. Storing JSON means adding a new field to
\code{SessionState} requires no database migration; the code just works. The price is that
you cannot query on those fields (``show me all calls stuck on the phone slot'' is not
expressible in SQL here), and nothing validates the shape. For a single-purpose store that is
only ever read back whole, by primary key, this is a sound trade. It would not be for a table
you needed to report on.
\end{keyidea}

The \code{save} method uses an \emph{upsert}, which inserts a row or updates it if the key
already exists, in one atomic statement:

\begin{lstlisting}[language=SQL,caption={The upsert in session\_store.save}]
INSERT INTO sessions (call_sid, state_json, created_at, updated_at)
VALUES (?, ?, ?, ?)
ON CONFLICT(call_sid) DO UPDATE
  SET state_json = excluded.state_json, updated_at = excluded.updated_at
\end{lstlisting}

\begin{jargon}{Parameterised query (the \code{?} marks)}
The \code{?} placeholders are not string formatting. The database receives the query and the
values separately, so a value can never be mistaken for SQL code. This is the correct defence
against \emph{SQL injection}, where hostile input like \code{'; DROP TABLE appointments; --}
would otherwise be executed as a command. Every query in this project is parameterised, which
is exactly right, and worth pointing out because a caller's spoken name goes into the database.
\end{jargon}

\code{mark\_abandoned\_if\_unresolved} is the neat part. When Twilio reports the call ended,
this checks whether any outcome was ever reached; if not, the call is recorded as
\code{abandoned}. Crucially it refuses to overwrite an outcome that already exists, so hanging
up after a successful booking still counts as booked. There is a test for exactly this.

\subsection{\file{voice\_agent/calendar\_store.py}}

The largest module (339 lines) and the one that does the real work. It defines the
\code{CalendarStore} interface, two implementations, and the conflict rule.

\subsubsection{The conflict rule}

Everything hinges on eight characters:

\begin{lstlisting}[language=Python,caption={The half-open interval overlap test}]
def _overlaps(a_start, a_end, b_start, b_end) -> bool:
    """Half-open interval overlap: touching endpoints (a_end == b_start) is NOT a conflict."""
    return a_start < b_end and b_start < a_end
\end{lstlisting}

\begin{keyidea}{Why \code{<} and not \code{<=}}
Two appointments where one ends at exactly 9:30 and the next starts at exactly 9:30 do not
overlap; they are back to back, which is what a busy clinic wants. Using \code{<=} would
report a false conflict and refuse a perfectly good booking. This is the classic
\emph{half-open interval}: the start instant belongs to the appointment, the end instant does
not.
\end{keyidea}

\begin{figure}[H]
\centering
\resizebox{0.95\textwidth}{!}{
\begin{tikzpicture}[xscale=1.5]
  \draw[->, thick, inkblue] (0,0) -- (9.3,0) node[right, font=\small] {time};
  \foreach \x/\lab in {0.5/9:00, 2.5/9:30, 4.5/10:00, 6.5/10:30, 8.5/11:00}
    { \draw[linegrey] (\x,-0.15) -- (\x,4.6); \node[font=\scriptsize, below] at (\x,-0.15) {\lab}; }

  % existing
  \fill[accent!30, draw=accent, thick] (0.5,3.7) rectangle (2.5,4.3);
  \node[font=\scriptsize, inkblue] at (1.5,4.0) {EXISTING 9:00-9:30};

  % case 1 adjacent
  \fill[okgreen!25, draw=okgreen, thick] (2.5,2.8) rectangle (4.5,3.4);
  \node[font=\scriptsize] at (3.5,3.1) {9:30-10:00};
  \node[font=\scriptsize\bfseries, okgreen, right] at (4.7,3.1) {NO conflict: touches only at the boundary};

  % case 2 partial
  \fill[warnred!25, draw=warnred, thick] (1.5,1.9) rectangle (3.5,2.5);
  \node[font=\scriptsize] at (2.5,2.2) {9:15-9:45};
  \node[font=\scriptsize\bfseries, warnred, right] at (4.7,2.2) {CONFLICT: partial overlap};

  % case 3 exact
  \fill[warnred!25, draw=warnred, thick] (0.5,1.0) rectangle (2.5,1.6);
  \node[font=\scriptsize] at (1.5,1.3) {9:00-9:30};
  \node[font=\scriptsize\bfseries, warnred, right] at (4.7,1.3) {CONFLICT: exact same slot};

  % case 4 later
  \fill[okgreen!25, draw=okgreen, thick] (4.5,0.1) rectangle (6.5,0.7);
  \node[font=\scriptsize] at (5.5,0.4) {10:00-10:30};
  \node[font=\scriptsize\bfseries, okgreen, right] at (6.7,0.4) {NO conflict};
\end{tikzpicture}}
\caption{\code{a\_start < b\_end and b\_start < a\_end}, drawn. The green 9:30 case is the one
a naive \code{<=} implementation gets wrong, and it is covered by scenario
\code{boundary\_adjacent\_appointments\_no\_conflict}.}
\end{figure}

\subsubsection{\code{SqliteCalendarStore}, the default}

One table, \code{appointments}, with \code{status} either \code{booked} or \code{cancelled}.
Cancelling is a \emph{soft delete}: the row stays, the status changes, so the history survives
and the conflict check simply ignores cancelled rows.

\begin{jargon}{Soft delete}
Marking a record as inactive rather than physically removing it. You keep an audit trail and
can undo; the cost is that every query must remember to filter out the dead rows.
\end{jargon}

\begin{honest}{\code{find\_conflict} reads the entire table into Python on every booking}
The query is \code{SELECT * FROM appointments WHERE status = 'booked'} with \emph{no time
range filter at all}. Every booked appointment ever made, past and future, is loaded and
looped over in Python to find an overlap. There is also no index on \code{start\_time} or
\code{phone}.

At this project's scale (one phone line, a handful of appointments) this is completely
invisible, and choosing simplicity over premature optimisation is defensible. But it is
\emph{O(n)} work per booking that a database is built to do in \emph{O(log n)}, and the fix is
a one-line \code{WHERE start\_time < ? AND ...} clause plus an index. Worth being able to say
``I know, here is exactly why it does not matter yet, and here is the fix'' rather than being
shown it.
\end{honest}

\begin{jargon}{O(n) and O(log n)}
Shorthand for how work grows with data size. \emph{O(n)}: ten times the appointments means ten
times the work. \emph{O(log n)}: ten times the appointments means only slightly more work,
because an index lets the database skip most rows. This is what indexes are for.
\end{jargon}

\subsubsection{\code{GoogleCalendarStore}, the optional adapter}

\begin{sloppypar}
A complete second implementation against the real Google Calendar API v3, using
\code{events.insert}, \code{events.patch}, \code{events.delete} and \code{freebusy.query} for
conflict detection. \code{get\_calendar\_store()} selects it automatically when the environment
variable \code{GOOGLE\_CALENDAR\_CREDENTIALS\_PATH} is set and the file it names exists;
otherwise the SQLite store is used.
\end{sloppypar}

\begin{jargon}{Adapter (design pattern)}
Code whose only job is to make one thing present itself with the shape of another. Here it
makes Google Calendar look like a \code{CalendarStore}, translating between appointment ids and
Google event ids so nothing upstream has to know the difference.
\end{jargon}

The Google client libraries are imported \emph{inside} \code{\_\_init\_\_}, not at the top of
the file. That is deliberate: the default SQLite install then does not need them present.

\begin{honest}{The Google adapter has never been run, and it has at least one real defect}
The README is upfront that no Google Calendar credentials were ever provided, so this class
has never executed against a live calendar. It is not a stub, but it is unverified. Reading it
closely, two concrete problems are visible without running it:

\textbf{One.} \code{cancel()} calls \code{events().delete()}, which destroys the event
outright, then \emph{fabricates} a return value with \code{start\_time=datetime.now()}. But
\code{conversation.py} uses that return value to speak to the caller:
\code{f"Done, your appointment for \{when\} is cancelled"}. On the Google backend the agent
would therefore read back \emph{the current time} rather than the appointment's real time. On
SQLite this is invisible, because there \code{cancel()} returns the true row.

\textbf{Two.} The \code{Appointment} dataclass declares \code{id: int}, but Google event ids
are strings, and this class assigns them straight into that field. Python does not enforce
type hints at runtime so nothing crashes, but the declared contract is being broken, and
\code{SessionState.target\_appointment\_id} is likewise typed \code{Optional[int]}.

Neither is hard to fix. Both are the predictable cost of writing an implementation you cannot
run, which is itself the honest lesson here.
\end{honest}

\subsection{\file{voice\_agent/intent\_engine.py}}

This is where the language model lives. \code{TurnContext} in, \code{TurnResult} out.

\begin{lstlisting}[language=Python,caption={The data crossing the boundary}]
@dataclass
class TurnContext:
    caller_number: str
    utterance: str                      # "" if <Gather> captured no speech
    history: list[dict]                 # [{"speaker": "agent"/"caller", "text": ...}]
    known_slots: dict                   # what we have collected so far
    current_action: Optional[str]       # book/reschedule/cancel, if established
    today: date                         # injected, never datetime.now()

@dataclass
class TurnResult:
    reply: str                          # what to say back
    intent: str                         # book|reschedule|cancel|unclear|goodbye
    name: Optional[str] = None
    phone: Optional[str] = None
    date_iso: Optional[str] = None      # "YYYY-MM-DD"
    time_24h: Optional[str] = None      # "HH:MM"
    duration_minutes: Optional[int] = None
    confirmed: Optional[bool] = None
\end{lstlisting}

\begin{keyidea}{\code{confirmed} is three-valued, and that is deliberate}
\code{Optional[bool]} means it can be \code{True}, \code{False}, or \code{None}, and all three
are distinct: ``yes, book it'', ``no, that's wrong'', and ``I have not been asked yet''. A plain
\code{True/False} would silently merge ``said no'' with ``never asked'', and the agent would
either book without consent or never book at all. The same three-valued logic is why the code
tests \code{if slots.get("confirmed") is not True} rather than \code{if not confirmed}.
\end{keyidea}

\subsubsection{Making the model fast enough for a phone call}

Three settings do the work, and the reasoning behind them is the most interesting engineering
judgement in the project:

\begin{lstlisting}[language=Python,caption={GeminiIntentEngine.parse\_turn, the configuration}]
config=types.GenerateContentConfig(
    system_instruction=system_prompt,
    thinking_config=types.ThinkingConfig(thinking_budget=0),  # turn thinking OFF
    response_mime_type="application/json",                    # must be JSON
    response_schema=self._response_schema(),                  # this exact shape
    temperature=0.1,                                          # near-deterministic
)
\end{lstlisting}

\begin{keyidea}{Sometimes engineering means turning the clever feature off}
Gemini 2.5 Flash ``thinks'' before answering by default, which improves hard reasoning and
costs seconds and tokens. The README records an actual measurement: an exploratory call spent
56 thinking tokens on a one-word answer. On a phone call, multiple seconds of dead air is a
product failure; the caller assumes the line dropped. So thinking is set to zero. The measured
result with thinking off plus a strict schema was 1.8 seconds and a clean structured reply.

This is a genuinely good instinct to be able to articulate: the right configuration depends
entirely on the product, and for a live phone turn, latency beats depth.
\end{keyidea}

\begin{jargon}{Token}
The unit language models read, write, and bill in: roughly a short word or word-fragment.
``Thinking tokens'' are ones the model generates privately to reason before answering. You pay
for them, and you wait for them.
\end{jargon}

\begin{jargon}{Temperature}
A dial for randomness. At 0 the model gives its single most likely answer every time; higher
values invite variety. For creative writing you want it high. For extracting a date from a
sentence you want it near zero, hence 0.1.
\end{jargon}

\begin{jargon}{Response schema, and structured output}
Rather than asking the model in prose to ``please reply in JSON'' and hoping, the schema is
supplied to the API as a formal specification and the platform constrains the output to match
it. The difference matters: politely asking sometimes yields ``Sure! Here's the JSON:
\{...\}'', which then fails to parse. This is the robust way to get machine-readable output.
\end{jargon}

\subsubsection{Defensive handling of the reply}

The parsing is properly paranoid, which is correct when the input comes from a model:

\begin{lstlisting}[language=Python,caption={Never trust the model's output shape}]
intent = data.get("intent", "unclear")
if intent not in VALID_INTENTS:
    intent = "unclear"           # an unknown intent degrades to "ask again"
...
reply=data.get("reply") or "Sorry, could you say that again?",
name=data.get("name") or None,   # "" and null both collapse to None
\end{lstlisting}

Any exception at all from the API call becomes an \code{IntentEngineError}, which
\code{conversation.py} turns into a graceful spoken handoff rather than a crash. An unknown
intent degrades to \code{unclear}, which routes into the retry path. Both are the right
failure direction: a caller hears a sentence, never silence.

\subsection{\file{voice\_agent/conversation.py}: the state machine}

The heart of the project. 338 lines, no Twilio, no Gemini, no SQL.

\begin{figure}[H]
\centering
\resizebox{0.93\textwidth}{!}{
\begin{tikzpicture}[node distance=7mm]
  \node[box] (start) {\code{handle\_turn(session, utterance)}};
  \node[dec, below=7mm of start, text width=20mm] (cap) {turn\_count\\ $>$ 14?};
  \node[box, right=22mm of cap, fill=warnred!12, draw=warnred] (h1) {\textbf{HANDOFF}\\{\scriptsize end call}};
  \node[box, below=7mm of cap] (parse) {build TurnContext,\\call \code{intent\_engine.parse\_turn}};
  \node[dec, below=7mm of parse, text width=20mm] (err) {engine\\raised?};
  % TECH FAILURE goes LEFT: the right column at this row is needed for the retry re-ask.
  \node[box, left=22mm of err, fill=warnred!12, draw=warnred] (h2) {\textbf{TECH FAILURE}\\{\scriptsize end call}};
  \node[dec, below=8mm of err, text width=20mm] (unc) {intent ==\\unclear?};
  \node[dec, right=24mm of unc, text width=20mm] (rc) {retry\_count\\ $>$ 3?};
  \node[box, right=20mm of rc, fill=warnred!12, draw=warnred] (h3) {\textbf{HANDOFF}};
  \node[box, above=9mm of rc, fill=softgrey] (rep) {``say that again?''\\{\scriptsize call continues}};
  \node[dec, below=8mm of unc, text width=20mm] (bye) {intent ==\\goodbye?};
  \node[box, right=22mm of bye, fill=okgreen!12, draw=okgreen] (end) {end call;\\{\scriptsize abandoned if no outcome yet}};
  \node[box, below=7mm of bye] (merge) {\code{retry\_count = 0};\\merge slots; set/switch action};
  \node[dec, below=7mm of merge, text width=22mm] (act) {which\\action?};
  \node[box, below left=9mm and 8mm of act] (bk) {\code{\_handle\_book}};
  \node[box, below=9mm of act] (cx) {\code{\_handle\_cancel}};
  \node[box, below right=9mm and 8mm of act] (rs) {\code{\_handle\_reschedule}};

  \draw[flow] (start) -- (cap);
  \draw[flow] (cap) -- node[lbl,above]{yes} (h1);
  \draw[flow] (cap) -- node[lbl,left]{no} (parse);
  \draw[flow] (parse) -- (err);
  \draw[flow] (err) -- node[lbl,above]{yes} (h2);
  \draw[flow] (err) -- node[lbl,right]{no} (unc);
  \draw[flow] (unc) -- node[lbl,above]{yes} (rc);
  \draw[flow] (rc) -- node[lbl,above]{yes} (h3);
  \draw[flow] (rc) -- node[lbl,right]{no} (rep);
  \draw[flow] (unc) -- node[lbl,left]{no} (bye);
  \draw[flow] (bye) -- node[lbl,above]{yes} (end);
  \draw[flow] (bye) -- node[lbl,left]{no} (merge);
  \draw[flow] (merge) -- (act);
  \draw[flow] (act) -- node[lbl,left]{book} (bk);
  \draw[flow] (act) -- node[lbl,left]{cancel} (cx);
  \draw[flow] (act) -- node[lbl,right]{reschedule} (rs);
\end{tikzpicture}}
\caption{\code{handle\_turn}, the top-level decision flow. Note the order: the safety caps are
checked \emph{before} anything expensive happens, so a runaway call cannot keep billing Gemini.}
\label{fig:sm}
\end{figure}

\subsubsection{The two safety nets, and why there are two}

\begin{keyidea}{A retry cap and a turn cap catch different failures}
\code{MAX\_RETRY\_COUNT = 3} counts \emph{consecutive unintelligible} turns. It catches the
bad line, the caller in a tunnel, the background noise.

\code{MAX\_TURN\_COUNT = 14} counts \emph{all} turns, and catches the case the retry cap
cannot: a caller whose every turn is perfectly intelligible but never actually completes
anything. Say ``sometime in the morning'' forever and \code{retry\_count} resets to zero every
single turn, because each turn was a valid continuation. Without the second cap, the call
never ends. Scenario 24 tests exactly this, and asserts \code{retry\_count == 0} to prove it
was the turn cap and not the retry cap that fired.

That is a genuinely thoughtful pair of tests: the assertion proves \emph{which} mechanism
saved the call, not merely that something did.
\end{keyidea}

Note also \code{session.retry\_count = 0} is reset on any clear turn, so the cap means ``three
in a row'', not ``three per call''. That is the humane choice.

\subsubsection{The booking path, traced}

\code{\_handle\_book} runs a strict gauntlet, in this order:

\begin{pseudocode}{\_handle\_book, in order}
\Step{1. missing = any of (name, phone, date\_iso, time\_24h) still empty?}
\Indent{if so: return the model's question, keep the call open}
\Step{2. slots["confirmed"] is not True?}
\Indent{if so: return the model's read-back, keep the call open}
\Step{3. parse date+time into a real datetime}
\Indent{unparseable: clear time, clear confirmed, re-ask}
\Step{4. \_validate\_business\_slot(dt)}
\Indent{in the past, or outside 9:00-17:00: clear time, clear confirmed, re-ask}
\Step{5. calendar\_store.book(...)}
\Indent{ConflictError: clear time, clear confirmed, offer to pick another time}
\Step{6. success: outcome = "booked", speak the confirmation}
\end{pseudocode}

\begin{keyidea}{The clear-time-and-confirmation reset is the subtle bit}
Every rejection path does the same two things:
\code{slots["time\_24h"] = None} and \code{slots["confirmed"] = None}. Miss either and the bug
is nasty. Leave the time set and the next turn re-validates the same bad time forever. Leave
\code{confirmed} at \code{True} and the caller's ``yes'' to the \emph{rejected} 7am silently
becomes a ``yes'' to whatever time they mention next, booking without real consent. Resetting
both forces a fresh, explicit confirmation of the corrected time. This pattern is repeated
identically in \code{\_handle\_reschedule}.
\end{keyidea}

\subsubsection{Validation, and the injectable clock}

\begin{lstlisting}[language=Python,caption={Why now\_fn exists}]
def __init__(self, intent_engine, calendar_store, session_store,
             now_fn: Callable[[], datetime] = datetime.now):   # <- injected clock
    ...

def _validate_business_slot(self, dt):
    now = self.now_fn()                       # not datetime.now()
    if dt <= now:
        return "That date and time has already passed. ..."
    if not (config.BUSINESS_HOURS_START <= dt.hour < config.BUSINESS_HOURS_END):
        return "That's outside our business hours ..."
    return None
\end{lstlisting}

\begin{keyidea}{Injecting the clock is what makes the tests possible at all}
A test asserting ``the appointment lands on 2026-07-15 at 14:00'' after the caller says
``July 15th at 2pm'' is only correct on days when that date is still in the future. Call
\code{datetime.now()} directly and the test suite quietly rots: it passes today and starts
failing on 16 July 2026 for reasons unrelated to any code change. Threading a \code{now\_fn}
through and pinning every scenario to \code{FIXED\_NOW = datetime(2026, 7, 6, 13, 0)} makes
time itself an input. Notice \code{FIXED\_NOW} is at 13:00: mid-afternoon, deliberately, so a
same-day ``past'' time (10am) can be tested without also tripping the business-hours check,
keeping the two rules independently testable. That is careful test design.
\end{keyidea}

Note the business-hours check tests \code{dt.hour} only, so 16:45 passes even though a 30
minute appointment then runs 15 minutes past closing. \code{config.py} concedes this in a
comment (``last bookable start is 16:xx for a 30 min slot''). It is a known simplification, not
an oversight.

\subsubsection{Changing your mind mid-call}

\begin{lstlisting}[language=Python,caption={The action switch}]
if session.action is None:
    session.action = result.intent
elif session.action != result.intent:
    # Caller changed their mind mid-call: start this new action clean.
    session.action = result.intent
    session.slots = {}
    session.target_appointment_id = None
\end{lstlisting}

Wiping the slots on a switch is right: a date gathered for a booking should not leak into a
cancellation. The cost is that a caller who switches and switches back re-answers everything.
For a phone agent, correctness beats convenience here.

\subsection{\file{voice\_agent/app.py}: the web layer}

Deliberately thin. Its whole job is: receive an HTTP form, call the manager, turn the answer
into TwiML.

\begin{jargon}{FastAPI}
A Python web framework: a library that handles the plumbing of receiving HTTP requests and
routing them to your functions, so you write only the logic. The \code{@app.post("/voice/gather")}
line above a function means ``run this when a POST arrives at that address''.
\end{jargon}

\subsubsection{The app factory, and a clever import trick}

\code{create\_app(...)} takes optional stores and an engine. Pass nothing and it wires up the
real production dependencies; pass fakes and you get a test instance. But that creates a
problem: \code{uvicorn} needs a module-level \code{app} object, and building one would demand
real credentials, which would make merely \emph{importing} the module fail in tests.

\begin{lstlisting}[language=Python,caption={PEP 562 module-level \_\_getattr\_\_}]
def __getattr__(name: str):
    if name == "app":
        return create_app()
    raise AttributeError(f"module {__name__!r} has no attribute {name!r}")
\end{lstlisting}

\begin{keyidea}{Lazy construction via module \code{\_\_getattr\_\_}}
This hook runs only when someone actually asks for an attribute the module does not otherwise
have. So \code{import voice\_agent.app} costs nothing and needs no credentials, while
\code{uvicorn voice\_agent.app:app} touches \code{app}, triggers the hook, and builds the real
thing. It neatly resolves the tension between ``uvicorn wants a module-level object'' and
``tests must import this without credentials''. It is niche but legitimate Python (PEP 562),
and it is well commented in the source.
\end{keyidea}

\subsubsection{Signature verification}

\begin{jargon}{Webhook signature}
Your webhook address is a public URL: anyone who finds it can POST to it and pretend to be
Twilio. To prevent this, Twilio signs each request with a cryptographic hash computed from the
URL, the form data, and your secret auth token, in the \code{X-Twilio-Signature} header. You
recompute it with the same token; if they match, the request really came from Twilio, because
only Twilio and you know the token.
\end{jargon}

\begin{honest}{The signature check silently does nothing in the default configuration}
This is the most significant finding in the review, and it was confirmed by running the code,
not by reading it.

\begin{lstlisting}[language=Python]
if app.state.enforce_signature:
    base_url = config.public_webhook_base_url()
    auth_token = os.environ.get("TWILIO_AUTH_TOKEN", "")
    if base_url and auth_token:              # <- if either is missing, NO CHECK RUNS
        if not verify_twilio_signature(request, form, auth_token, base_url):
            return Response(status_code=403, content="Invalid Twilio signature")
\end{lstlisting}

\code{PUBLIC\_WEBHOOK\_BASE\_URL} is blank by default, and \file{.env.example} explicitly
instructs ``Leave blank for local development and testing''. So with \code{enforce\_signature=True},
the nominally secure setting, an entirely unsigned request is accepted. Driving the real app
with a \code{TestClient} confirms it: base URL unset plus no signature header at all returns
\textbf{HTTP 200}. Set the base URL and the same unsigned request correctly returns
\textbf{HTTP 403}, so the check itself is implemented correctly.

There is a real reason for the design: Twilio signs the \emph{public} URL it requested, so
without knowing that URL you cannot recompute the signature, and there is genuinely nothing to
check against locally. The danger is the failure \emph{direction}. This \emph{fails open}: the
consequence of a missing setting is silently no security, not a refusal to start. Deploy this
behind a real URL and forget the variable, and the booking endpoints are wide open to anyone
who finds them, with no warning in the logs.

The fix is small and worth stating in an interview: if \code{enforce\_signature} is true and
the base URL is missing, refuse to start, or at minimum log a loud warning on every request.
Make the safe path the default and the insecure path something you must deliberately choose,
for example an explicit \code{allow\_unsigned=True}.
\end{honest}

\subsubsection{Endpoint behaviour}

\code{/voice/status} accepts Twilio's terminal statuses (\code{completed}, \code{failed},
\code{busy}, \code{no-answer}, \code{canceled}) and marks the call abandoned if no outcome was
reached. It returns \code{204 No Content}, which is correct: Twilio wants nothing back.

\subsection{\file{voice\_agent/cli.py}: the outbound test call}

The only code in the repository that spends real money on telephony, and it is fenced off
accordingly. In order: validate the number against an E.164 regex; refuse outright without a
public URL; load and validate the Twilio credentials; then an interactive
\code{[y/N]} confirmation naming both numbers and warning of charges, skippable only with an
explicit \code{--yes}.

\begin{honest}{The guardrails are real, and this document respected them}
Nothing else in the repository imports or invokes \code{cli.py}; it runs only when a human
types the command. In line with the project's budget rule and the instructions for this
review, \textbf{no phone call was placed and no billed API test was run} while producing this
document. The 7 tests marked \code{real\_api} were deliberately deselected, never executed.
Every claim below about them rests on reading the code and the README, and is labelled as such.
\end{honest}

\section{The evaluation harness}

\subsection{What it actually is}

\code{eval/scenarios.py} defines 25 \code{Scenario} objects. Each is one or more simulated
calls (a phone number plus a list of caller utterances) and a \code{check} function asserting
on the resulting calendar and session state.

\begin{keyidea}{What the 25 scenarios do and do not exercise}
This is the claim to pin down precisely, because the honest answer is narrower than
``25 scenarios pass'' sounds, and the repository states it correctly.

\textbf{Real, and genuinely exercised:} the \code{ConversationManager} state machine, the real
\code{SqliteCalendarStore} on a real temporary database file, the real \code{SessionStore},
real conflict detection, real business-hours and past-date validation, real retry and turn
caps.

\textbf{Replaced by a test double:} the language model, swapped for the regex-based
\code{ScriptedIntentEngine}.

\textbf{Not involved at all:} Twilio, FastAPI, HTTP, TwiML. Utterances are fed straight into
\code{manager.handle\_turn}, deliberately ``bypassing Twilio and FastAPI entirely'' as the
module docstring says.

So ``25/25'' is a strong statement about the booking \emph{logic}, and no statement whatever
about speech recognition or the phone line. The FastAPI and TwiML layers are covered
separately, by the 7 endpoint tests. The README does not overstate this.
\end{keyidea}

\begin{figure}[H]
\centering
\resizebox{0.9\textwidth}{!}{
\begin{tikzpicture}[node distance=8mm]
  \node[box] (s) {\code{run\_scenario(scenario)}};
  \node[box, below=6mm of s] (tmp) {\code{tempfile.mkdtemp()}\\{\scriptsize a fresh throwaway directory}};
  \node[box, below=6mm of tmp] (wire) {build real SqliteCalendarStore + real SessionStore\\+ ScriptedIntentEngine + \code{now\_fn = lambda: FIXED\_NOW}};
  \node[box, below=6mm of wire] (loop) {for each call, for each utterance:\\\code{manager.handle\_turn(session, text)}\\{\scriptsize break early if outcome.end\_call}};
  \node[box, below=6mm of loop] (chk) {\code{scenario.check(result)}\\{\scriptsize assert on calendar + session state}};
  \node[box, below=6mm of chk, fill=warnred!8, draw=warnred] (fin) {\code{finally: shutil.rmtree(tmp\_dir)}\\{\scriptsize total isolation, always cleaned up}};
  \draw[flow] (s)--(tmp); \draw[flow] (tmp)--(wire); \draw[flow] (wire)--(loop);
  \draw[flow] (loop)--(chk); \draw[flow] (chk)--(fin);
\end{tikzpicture}}
\caption{Every scenario gets a brand new temporary database and is destroyed afterwards in a
\code{finally} block, so no scenario can ever influence another.}
\end{figure}

A nice defensive touch: \code{scenario.check} is itself wrapped in a \code{try/except}, so a
bug in a check function fails that one scenario rather than crashing the whole run.

\subsection{The 25 scenarios, verified}

The harness was run for this document. It required no third-party packages and made no network
calls, which independently confirms the ``runs free'' claim.

\begin{lstlisting}[caption={Observed output, reproducing the committed docs/screenshots/run-output.txt exactly}]
$ python -m eval.run_eval
[PASS] (book) book_simple_one_shot
[PASS] (book) book_multi_turn_slot_filling
... 21 lines omitted for space ...
[PASS] (unclear_speech) max_turn_cap_handoff
[PASS] (conflict) reschedule_conflict_then_alternate

25/25 scenarios passed
\end{lstlisting}

\begin{table}[H]
\centering
\small
\begin{tabular}{@{}rlp{7.7cm}@{}}
\toprule
\textbf{\#} & \textbf{Category} & \textbf{What it actually pins down} \\
\midrule
1 & book & Everything in one sentence, then ``yes''. The happy path. \\
2 & book & One fact per turn; a spoken phone number wins over caller ID. \\
3 & book & ``one hour'' overrides the 30 minute default. \\
4 & book & Spoken number stored; caller ID \emph{not} used. See finding below. \\
5 & book & Caller rejects the read-back, corrects, confirms: exactly 1 appointment, not 2. \\
6 & reschedule & Books in call 1, moves it in call 2. State survives across calls. \\
7 & reschedule & Nothing on file: graceful, then goodbye becomes \code{abandoned}. \\
8 & cancel & Cancel marks status \code{cancelled}, active list empties, row survives. \\
9 & cancel & Cancel with nothing on file. \\
10 & conflict & Second caller wants a taken slot, is refused, books another. \\
11 & ambiguous & ``sometime in the morning'' re-asks rather than guessing. \\
12 & ambiguous & ``later this week'' re-asks; concrete date then accepted. \\
13 & unclear & Silence once, recovers; \code{retry\_count} back to 0. \\
14 & unclear & Four unusable turns: retry cap fires, handoff. \\
15 & abandoned & Twilio status callback marks an incomplete call abandoned. \\
16 & book & 7am confirmed, rejected as too early, corrected to 9am. \\
17 & book & 9pm confirmed, rejected as too late, corrected to 4pm. \\
18 & book & Today 10am (past, vs \code{FIXED\_NOW} 13:00), rejected, corrected. \\
19 & conflict & \textbf{9:00 and 9:30 both book.} The half-open boundary. \\
20 & conflict & 10:15 overlaps a 10:00 slot; refused; 11:00 accepted. \\
21 & book & Mid-call switch from book to cancel; slots wiped. \\
22 & book & Two callers, two slots, independent. \\
23 & book & Goodbye after success does \emph{not} overwrite \code{booked}. \\
24 & unclear & Turn cap fires with \code{retry\_count == 0}. \\
25 & conflict & Reschedule into a taken slot; refused; alternate accepted. \\
\bottomrule
\end{tabular}
\caption{The 25 scenarios. 16, 17 and 18 each contain a deliberate extra ``yes'', because
validation only runs after confirmation; without it the scenario would pass without ever
executing the code it is named for.}
\end{table}

\subsection{Do the scenarios have teeth? A verification of the verification}

A passing test proves nothing unless it can fail. The README claims the caps were deliberately
broken to check the tests noticed. That claim was independently re-tested for this document, by
copying the project to a scratch directory and sabotaging it there (the repository itself was
never modified):

\begin{table}[H]
\centering
\small
\begin{tabular}{@{}lll@{}}
\toprule
\textbf{Sabotage applied} & \textbf{Result} & \textbf{Scenario that caught it} \\
\midrule
\code{MAX\_RETRY\_COUNT} 3 $\rightarrow$ 999 & 24/25 & \code{unclear\_speech\_retry\_cap\_handoff} \\
\code{MAX\_TURN\_COUNT} 14 $\rightarrow$ 9999 & 24/25 & \code{max\_turn\_cap\_handoff} \\
\code{<} $\rightarrow$ \code{<=} in \code{\_overlaps} & 24/25 & \code{boundary\_adjacent\_appointments\_no\_conflict} \\
\bottomrule
\end{tabular}
\caption{Mutation testing, performed for this document. Each sabotage killed exactly the one
scenario named for it, and no others. The scenarios are load-bearing, not decorative.}
\end{table}

\begin{jargon}{Mutation testing}
Deliberately introducing a bug to check that the test suite fails. It tests the tests. If you
break the code and everything still passes, your tests were measuring nothing. The precision
here (each mutation killed exactly one scenario, no collateral) is a sign of well-targeted
tests.
\end{jargon}

\subsection{The scripted engine}

Roughly a dozen regexes over keywords: \code{\_BOOK\_KW}, \code{\_CANCEL\_KW},
\code{\_YES\_KW}, plus extractors for names, phones, times, dates and durations. Its
resolution order matters:

\begin{lstlisting}[language=Python,caption={Intent resolution, in priority order}]
if _CANCEL_KW.search(utterance):        new_intent = "cancel"
elif _RESCHEDULE_KW.search(utterance):  new_intent = "reschedule"
elif _BOOK_KW.search(utterance):        new_intent = "book"
elif _GOODBYE_KW.search(utterance):     new_intent = "goodbye"
...
if new_intent is not None:                     intent = new_intent
elif ctx.current_action is not None and (extracted_anything or vague_answer):
    intent = ctx.current_action    # a continuation of what we were doing
else:
    intent = "unclear"             # drives the retry path
\end{lstlisting}

The \code{vague\_answer} branch is the clever part. ``Sometime in the morning'' extracts no
usable slot, but it is not gibberish either, and conflating the two would make scenario 24
untestable: a real answer that happens to be unhelpful must re-ask \emph{without} incrementing
the retry counter.

\begin{honest}{The scripted engine flatters the state machine, unavoidably}
This is worth naming clearly because it bounds what ``25/25'' means. The scripted engine was
written \emph{after} the scenarios, to understand exactly the phrases those scenarios use. Its
own docstring admits this: ``its natural-language coverage only needs to be as good as the
phrases the scenarios use''. So the 25 scenarios can never surface a misunderstanding, because
the fake model never misunderstands anything; the sentences and the parser were built for each
other.

That is the correct trade (the alternative is 100+ billed calls per run, and flaky tests), and
the project is open about it. But it means the scenarios test the state machine \emph{given
perfect understanding}. Real understanding is tested only by the 6 marked Gemini tests, against
clean typed English. Nothing here tests messy speech-to-text output. The right framing under
questioning: ``25/25 proves the booking logic is correct, and I deliberately made the
understanding layer a separate, smaller, honest experiment.''
\end{honest}

Note too that the parsers are duplicated in spirit: \code{ScriptedIntentEngine} reimplements
relative-date logic (``tomorrow'', ``next Tuesday'') that the Gemini prompt asks the model to
do. Two implementations of one rule can drift, though only the Gemini one ships.

\section{Testing: what was verified, and how}

\subsection{The reported numbers, independently reproduced}

Every count below was re-verified for this document by installing the pinned requirements into
a brand new virtual environment and running the suite. The free suite was run in full; the
billed suite was collected but \textbf{deliberately never executed}.

\begin{table}[H]
\centering
\small
\begin{tabular}{@{}lll@{}}
\toprule
\textbf{Claim in the README} & \textbf{Observed} & \textbf{Verdict} \\
\midrule
54 tests collected & \code{54 tests collected in 0.72s} & confirmed \\
47 run free, 7 billed & \code{47 passed, 7 deselected} & confirmed \\
7 marked \code{real\_api} & \code{7/54 tests collected (47 deselected)} & confirmed, \emph{not run} \\
25/25 scenarios pass & \code{25/25 scenarios passed} & confirmed \\
Installs into a clean venv & installed, exit code 0 & confirmed \\
6 real Gemini calls & 6 test functions, 1 call each & \emph{inferred from source} \\
1.8 s with thinking off & no artefact in the repository & \emph{unverifiable here} \\
Live tunnel booking & no artefact in the repository & \emph{unverifiable here} \\
\bottomrule
\end{tabular}
\caption{Verification of the README's own claims. The free numbers all hold exactly. The
last three are the owner's own records of runs that left no committed evidence.}
\label{tab:verify}
\end{table}

\begin{honest}{Three claims rest on the owner's word, not on anything in the repository}
Stated plainly so they are not defended as if verified:

\textbf{The 1.8 second latency measurement} and the ``56 thinking tokens'' figure appear only
in README prose. No script, log, or artefact reproduces them. They are plausible and specific,
but a sceptical interviewer cannot check them, and neither could this review.

\textbf{The live public-tunnel run} (the \code{localhost.run} tunnel, the genuinely
HMAC-signed requests, the row written to the database and then deleted) is the strongest
verification claimed, and nothing of it survives in the repository: no log, no transcript, no
recording. It was, by design, cleaned up afterwards.

\textbf{The Gemini cost figure} is explicitly labelled in the README as ``an estimate from
per-token rates, not a verified line-item invoice'', which is exactly the right way to say it.

None of this suggests the claims are untrue; the README's overall precision argues strongly the
other way, and its willingness to flag its own estimate as an estimate is a good sign. But the
distinction between ``I ran this and here is the artefact'' and ``I ran this and cleaned up''
is one an interviewer may probe. The cheap fix, for next time, is to commit the log.
\end{honest}

\subsection{The coverage map}

\begin{figure}[H]
\centering
\resizebox{0.95\textwidth}{!}{
\begin{tikzpicture}[node distance=5mm]
  \node[box, fill=okgreen!15, draw=okgreen, text width=34mm] (a) {\textbf{calendar\_store.py}\\{\scriptsize 13 tests, no network}\\{\scriptsize + 25 scenarios}};
  \node[box, fill=okgreen!15, draw=okgreen, text width=34mm, right=5mm of a] (b) {\textbf{conversation.py}\\{\scriptsize 25 scenarios}\\{\scriptsize + endpoint tests}};
  \node[box, fill=okgreen!15, draw=okgreen, text width=34mm, right=5mm of b] (c) {\textbf{session\_store.py}\\{\scriptsize exercised throughout}};
  \node[box, fill=okgreen!15, draw=okgreen, text width=34mm, right=5mm of c] (d) {\textbf{app.py / twiml.py}\\{\scriptsize 7 TestClient tests}};

  \node[box, fill=accent!15, draw=accent, text width=34mm, below=7mm of a] (e) {\textbf{GeminiIntentEngine}\\{\scriptsize 6 REAL billed tests,}\\{\scriptsize clean typed English only}};
  \node[box, fill=accent!15, draw=accent, text width=34mm, below=7mm of b] (f) {\textbf{Twilio signature}\\{\scriptsize real token, local HMAC,}\\{\scriptsize no network}};
  \node[box, fill=warnred!12, draw=warnred, text width=34mm, below=7mm of c] (g) {\textbf{GoogleCalendarStore}\\{\scriptsize NEVER RUN}\\{\scriptsize no credentials provided}};
  \node[box, fill=warnred!12, draw=warnred, text width=34mm, below=7mm of d] (h) {\textbf{A real human voice call}\\{\scriptsize NEVER DONE}\\{\scriptsize needs a live tunnel + phone}};

  \node[lbl, left=3mm of a, text width=17mm] {\textbf{free +}\\\textbf{thorough}};
  \node[lbl, left=3mm of e, text width=17mm] {\textbf{narrow but}\\\textbf{real}};
\end{tikzpicture}}
\caption{Green: covered free and well. Blue: covered by a small, deliberate, real-API slice.
Red: not covered, and the README says so in both cases.}
\end{figure}

\section{Honest findings}
\label{sec:findings}

Everything below was found by reading the source, and each was confirmed by executing the real
code in a scratch copy. None of it appears in the README.

\begin{honest}{Finding 1: the wrong appointment is cancelled when a caller has two}
\code{\_find\_target\_appointment} takes the first match:

\begin{lstlisting}[language=Python]
matches = self.calendar_store.list_by_phone(phone, active_only=True)
return matches[0] if matches else None
\end{lstlisting}

and \code{list\_by\_phone} sorts \code{ORDER BY start\_time DESC}, meaning \emph{furthest
future first}. So a caller with an appointment tomorrow and another next month who says
``cancel my appointment'' gets \textbf{next month's cancelled}, and tomorrow's silently
retained. Verified directly: booking 2026-08-01 and 2026-09-01 under one number and asking for
the target returns the September one.

The agent does read the date back (``I found your appointment for Tuesday September 1st,
should I cancel it?''), so an attentive caller can say no. But the default is wrong: ``my
appointment'' in ordinary speech means the next one. No scenario covers a caller with two
active appointments, which is exactly why the suite is green. The fix is \code{ASC} ordering
plus a filter to future appointments; the deeper fix is asking which one.
\end{honest}

\begin{honest}{Finding 2: phone numbers are stored in two incompatible formats}
Twilio supplies caller ID in E.164 (\code{+15551112222}). The scripted engine strips speech to
bare digits (\code{5551112222}). \code{conversation.py} fills the slot from whichever it has:

\begin{lstlisting}[language=Python]
if not session.slots.get("phone"):
    session.slots["phone"] = session.caller_number    # E.164 from Twilio
\end{lstlisting}

and \code{list\_by\_phone} does an exact string match. So an appointment booked after the
caller \emph{spoke} their number is stored under \code{5551112222} and is \textbf{unfindable}
when that same person phones back, because lookup then uses \code{+15551112222}. Verified:
booking under the spoken form and looking up the E.164 form returns 0 matches. Cancelling by
phone would fail for exactly the caller who was most helpful.

The sharpest part: scenario 4 (\code{book\_explicit\_phone\_override}) \emph{asserts this
behaviour as correct}, checking \code{len(by\_caller\_id) == 0}. The test locks the bug in. Its
intent (a spoken number should win over caller ID) is right; what is missing is normalising
both to one canonical format on the way in. This is a good illustration of a test that
encodes an assumption rather than a requirement.
\end{honest}

\begin{honest}{Finding 3: cancelling an already-cancelled appointment silently succeeds}
The \code{CalendarStore} interface says \code{cancel} ``Raises NotFoundError''.
\code{SqliteCalendarStore.reschedule} does check status:

\begin{lstlisting}[language=Python]
existing = self.get(appointment_id)
if existing.status != "booked":
    raise NotFoundError(f"Appointment {appointment_id} is not active ...")
\end{lstlisting}

but \code{cancel} performs no such check and simply re-runs the \code{UPDATE}. Verified:
cancelling twice returns \code{status='cancelled'} both times with no error. The effect is
benign (the operation is idempotent, and arguably that is even desirable), but it is an
inconsistency between two sibling methods and a divergence from the documented contract. Worth
knowing which way you would resolve it.
\end{honest}

\begin{honest}{Finding 4: small dead code}
\begin{itemize}
  \item \code{SessionState.stage} is declared with a comment describing a lifecycle
  (\code{start -> collecting -> confirming -> done}) and is \textbf{never read or written
  anywhere}. It is serialised into every session row regardless. A grep for \code{stage} across
  the codebase returns exactly one hit: the declaration. It is a leftover from an earlier
  design that the \code{action} plus \code{slots} pair replaced.
  \item \code{\_handle\_cancel(self, session, model\_reply)} never uses \code{model\_reply};
  every return path writes its own sentence. Harmless, but it implies a symmetry with
  \code{\_handle\_book} that does not exist.
  \item \code{TurnOutcome.outcome} is populated carefully on every terminal path and
  \code{app.py} never reads it; only the tests and the eval harness do, and via
  \code{session.outcome} rather than the returned object.
  \item \code{DEEPGRAM\_API\_KEY} sits in \file{.env.example} and appears nowhere in the code.
  The file says so explicitly (``not wired into the code yet''), which makes it honest rather
  than misleading, but it is still an unused knob.
\end{itemize}
\end{honest}

\begin{honest}{Finding 5: \code{strftime("\%-I:\%M \%p")} is not portable}
Every caller-facing date uses the \code{\%-d} / \code{\%-I} form to strip a leading zero
(``at 2:00 PM'' rather than ``at 02:00 PM''). That hyphen flag is a glibc extension. It works
on Linux and macOS and raises \code{ValueError} on Windows. This project is very unlikely to
run on Windows, so it is a note rather than a defect, but it is the kind of thing that turns a
cosmetic detail into a crash in a caller-facing sentence.
\end{honest}

\begin{honest}{Finding 6: SQLite concurrency, acknowledged but worth being precise about}
The README correctly flags SQLite as single-writer. Two specifics sharpen it. Both stores open
their connection with \code{check\_same\_thread=False}, which switches off Python's own
safety check that a connection is not shared across threads; FastAPI serves requests from a
thread pool, so concurrent requests genuinely can share that connection object.

More interestingly, \code{book()} does \code{find\_conflict()} and then \code{INSERT} as two
separate statements with no transaction around them. Two simultaneous calls could both find no
conflict and both insert: a double booking, despite the conflict check. One phone line makes
this vanishingly unlikely, and it is the right call not to have solved it. But ``why is your
conflict check not atomic?'' is a fair question, and the answer (wrap both in a single
transaction, or add a database-level constraint) is worth having ready.
\end{honest}

\begin{jargon}{Race condition}
A bug where the result depends on the accidental timing of two things happening at once.
Finding 6 is the classic ``check then act'' race: the fact you checked (no conflict) can stop
being true between checking and acting on it.
\end{jargon}

\section{Security review}

\begin{table}[H]
\centering
\small
\begin{tabular}{@{}lp{9.4cm}@{}}
\toprule
\textbf{Check} & \textbf{Result} \\
\midrule
Secrets in the working tree & \textbf{None.} \file{.env} is a symbolic link pointing outside
the repository entirely, to a private secrets directory. \\
Secrets in git history & \textbf{None.} All 46 objects across the full history were scanned for
Twilio SIDs, Google API keys, private key blocks and hard-coded tokens. Clean. \\
\file{.env.example} & Placeholders only (\code{ACyour-account-sid-here}). No real value. \\
\file{.gitignore} & Covers \file{.env}, \code{service-account*.json}, \code{*.pem},
\code{*.db}, \code{data/}. Appropriate. \\
SQL injection & Every query is parameterised. No string-built SQL anywhere. \\
Credential logging & \code{config.py} never logs values. The Gemini error path logs the
exception, not the credential. \\
Webhook authentication & \textbf{Fails open by default.} See the finding in the \code{app.py}
section: the single most important issue found. \\
\bottomrule
\end{tabular}
\caption{Security posture. Secret hygiene is genuinely good and follows the portfolio's own
rule that credentials live outside the folder.}
\end{table}

\section{Dependencies, and why each is there}

\begin{table}[H]
\centering
\small
\begin{tabular}{@{}llp{6.5cm}@{}}
\toprule
\textbf{Package} & \textbf{Pin} & \textbf{Role} \\
\midrule
\code{fastapi} & 0.139.0 & The web framework Twilio's webhooks hit. \\
\code{uvicorn[standard]} & 0.50.2 & The server that actually runs FastAPI. \\
\code{python-multipart} & 0.0.32 & Parses form-encoded bodies. Twilio posts forms, not JSON,
so without this \code{await request.form()} fails. \\
\code{twilio} & 9.10.9 & TwiML builders, \code{RequestValidator}, and the REST client. \\
\code{google-genai} & 2.10.0 & The Gemini SDK, via Vertex AI. \\
\code{google-auth} & 2.55.1 & Service-account authentication. \\
\code{google-api-python-client} & 2.198.0 & Only for the optional Google Calendar adapter. \\
\code{python-dotenv} & 1.2.2 & Loads \file{.env} into the environment. \\
\code{pytest} & 9.1.1 & The test runner. \\
\code{httpx} & 0.28.1 & Required by FastAPI's \code{TestClient}. \\
\bottomrule
\end{tabular}
\caption{All ten dependencies are pinned to exact versions, which is right for reproducibility.
Installation into a clean virtual environment was verified for this document.}
\end{table}

\begin{jargon}{Pinning}
Recording the exact version (\code{==0.139.0}) rather than a range. It means the code you tested
is the code that runs. The cost is that upgrades become deliberate work, which for a portfolio
project is the right trade.
\end{jargon}

Note that \code{google-api-python-client} is installed for everyone but used only by the
unverified Google Calendar adapter. Correspondingly, \code{calendar\_store.py} imports it
lazily inside \code{\_\_init\_\_}, so the code is at least consistent with treating it as
optional.

\section{What to be ready to say about this project}

\begin{keyidea}{The strongest things here}
\begin{enumerate}
  \item \textbf{The interface seam.} Two small abstract classes turn an untestable,
  expensive, non-deterministic system into one with 47 free deterministic tests plus 25 free
  scenarios. This is the single best thing in the project and the easiest to defend.
  \item \textbf{Turning the LLM's flagship feature off.} Recognising that thinking mode is
  wrong for a live phone turn, measuring it, and disabling it, is real product engineering
  rather than cargo-culted best practice.
  \item \textbf{The tests were tested.} The mutation testing claim in the README is true; it
  was reproduced here. Each sabotage killed exactly its own scenario.
  \item \textbf{Two safety caps for two different failures}, with a test that proves \emph{which}
  one fired.
  \item \textbf{Reading the vendor documentation properly.} The \code{actionOnEmptyResult}
  discovery is exactly the class of detail that separates working telephony from a demo.
  \item \textbf{The README's honesty.} It is more candid about its own gaps than most
  production documentation.
\end{enumerate}
\end{keyidea}

\begin{keyidea}{The things to raise before an interviewer finds them}
\begin{enumerate}
  \item \textbf{The signature check fails open} without \code{PUBLIC\_WEBHOOK\_BASE\_URL}.
  Know the fix: refuse to start, or make unsigned an explicit opt-in.
  \item \textbf{Two phone number formats}, and a test that locks the bug in. Normalise on the
  way in.
  \item \textbf{The furthest-future appointment is targeted}, not the soonest.
  \item \textbf{The scripted engine cannot fail to understand}, so ``25/25'' means the logic is
  right given perfect understanding. Say so first; it is a strength that you know it.
  \item \textbf{No timezones at all}, acknowledged in the README.
  \item \textbf{The Google adapter has never run}, and reading it reveals the \code{cancel()}
  return-value defect.
  \item \textbf{The live tunnel run left no artefact.} Be ready to describe it precisely.
\end{enumerate}
\end{keyidea}

\vfill
\begin{center}
\small\itshape
Built from the repository source. Every claim above is either quoted from the code, verified by
running it in an isolated scratch copy, or explicitly labelled as inferred or unverifiable.
No phone call was placed and no billed API test was executed in the making of this document.
\end{center}

\end{document}
